% 2.2 超尖点表示（底本 p.16--20）
\subsection{超尖点表示}

\subsubsection{诱导与紧诱导}

我们在稍一般的情形下工作：设 $G$ 是局部紧（可以是完全不连通）拓扑群，
$H<G$ 是闭子群，考虑它们的光滑复表示。

\begin{definition}
定义两个函子 $\Ind,\ \ind\colon\Repinf_H\to\Repinf_G$：
\begin{enumerate}[label=(\arabic*),leftmargin=2.2em]
\item $\Ind^G_H(W):=\bigl\{f\colon G\to W\ \big|\ f(hx)=hf(x)\ \forall h\in H,\,x\in G;
\ \text{且}\ f\ \text{一致局部常值}\bigr\}$，
其中最后一个条件指：存在开子群 $C\subset G$（可依赖于 $x$）使 $f(xC)=f(x)$ 对一切 $x$ 成立。
$G$ 通过右平移作用在 $\Ind(W)$ 上。
\item $\ind^G_H\subset\Ind^G_H$ 是子函子，
\[
\ind(W)=\{f\in\Ind(W)\mid f\ \text{模}\ H\ \text{紧支集}\}.
\]
\end{enumerate}
\end{definition}

\begin{remark}
\leavevmode
\begin{enumerate}[label=(\arabic*),leftmargin=2.2em]
\item 若 $G/H$ 紧，则 $\ind=\Ind$。
\item 设 $G/H$ 紧，则 $\ind=\Ind$ 把 $H$ 的容许表示送到 $G$ 的容许表示。
\item $\Ind$ 正合（$\Res$ 限制函子也正合）；若进一步 $H<G$ 是开子群，则 $\ind$ 也正合。
\end{enumerate}
\end{remark}

考虑限制函子 $\Res\colon\Repinf_G\to\Repinf_H$。

\begin{lemma}
\leavevmode
\begin{enumerate}[label=(\arabic*),leftmargin=2.2em]
\item （Frobenius 互反）\en{6} 有伴随对 $\Res\dashv\Ind$。
\item 进一步设 $H<G$ 是开子群（于是 $G/H$ 离散），则 $\ind\dashv\Res$。
\end{enumerate}
\end{lemma}

\begin{remark}
\leavevmode
\begin{enumerate}[label=(\arabic*),leftmargin=2.2em]
\item 函子 $\ind$、$\Res$、$\Ind$ 在某种意义上对应于 $f_!$、$f^{*}$、$f_{*}$，
      这可以严格化。[例如参见 Ngo 的笔记] 记住上面的伴随对时，这个类比很有用。
\item 设 $H<G$ 是开子群，则 $\ind$ 可以解释为 $\mathrm{coInd}$，即
      $\ind(W)=W\otimes_{\mathbb{C}[H]}\mathbb{C}[G]$。
\end{enumerate}
\end{remark}

现在考虑商 $\psi\colon G\to M=G/N$，其中 $N$ 是某个（闭）正规子群。我们仍有「$\Res$」函子：
\[
\Res=\psi^{*}\colon\Repinf_M\longrightarrow\Repinf_G,
\]
由与 $\psi$ 复合给出（有时称为\emph{膨胀函子}）。$\Res$ 的左伴随是取 $N$-余不变量子空间，
$\Res$ 的右伴随是取 $N$-不变量。

我认为下面这个约定与一些已有的约定不同：

\begin{definition}
设 $\psi\colon G\to M=G/N$ 如上，定义
\[
\psi_{!}\colon\Repinf_G\to\Repinf_M,\qquad \psi_{!}V=V/[N-1]V,
\]
即 $V$ 的 $N$-余不变量空间。
\end{definition}

于是如所期望，在此设定下我们有伴随对 $\psi_{!}\dashv\psi^{*}$。

\subsubsection{抛物诱导与 Jacquet 函子 I}

设 $P<G=\GL_n(K)$ 是抛物子群\en{7}，与 $P$ 相伴有幂幺根 $U_P$ 与 Levi 商 $L_P$，
它们适合正合列
\[
1\longrightarrow U_P\longrightarrow P\longrightarrow L_P\longrightarrow 1 .
\]
我们有 $G=\GL_n(K)$，故抛物子群都是标准抛物子群的共轭：
\[
P_{\underline n}=
\begin{pmatrix}
\GL_{n_1} & * & *\\
 & \GL_{n_2} & *\\
 & & \ddots
\end{pmatrix},
\qquad \underline n=(n_1,\dots,n_r),
\]
其中 $\underline n$ 是 $n$ 的一个有序分拆。在 $P_{\underline n}$ 内有
\[
U_{\underline n}=
\begin{pmatrix}
1_{n_1} & * & *\\
 & 1_{n_2} & *\\
 & & \ddots
\end{pmatrix},
\qquad
L_{\underline n}=
\begin{pmatrix}
\GL_{n_1} & & \\
 & \ddots & \\
 & & \GL_{n_r}
\end{pmatrix}.
\]
现在对任一标准抛物子群 $P$（$P<G$，$\phi\colon P\twoheadrightarrow L=P/U$），
由上面的讨论我们有函子
\[
\ind=\Ind\colon\Repinf_P\to\Repinf_G,\qquad
\phi_{!}\colon\Repinf_P\to\Repinf_L,
\]
以及伴随对
\[
\Res^G_P\dashv\ind^G_P=\Ind^G_P,\qquad \phi_{!}\dashv\phi^{*}.
\]
现在定义（朴素 / 未正规化的）抛物诱导 $\widetilde I^G_P$ 与 $\widetilde J^G_P$ 如下：

\begin{definition}
\leavevmode
\begin{enumerate}[label=(\arabic*),leftmargin=2.2em]
\item $\widetilde I^G_P\colon\Repinf_L\to\Repinf_G$ 定义为 $\widetilde I^G_P:=\Ind^G_P\circ\phi^{*}$；
\item $\widetilde J^G_P\colon\Repinf_G\to\Repinf_L$ 定义为
      $\widetilde J^G_P=\phi_{!}\circ\Res^G_P$。
\end{enumerate}
\end{definition}

\begin{remark}
我们有以下性质：
\begin{enumerate}[label=(\arabic*),leftmargin=2.2em]
\item $\widetilde J\dashv\widetilde I$；
\item 两个函子都正合；
\item 两个函子都把有限长表示送到有限长表示；
\item 两个函子都把容许表示送到容许表示；
\item 上述性质在用特征标扭转之后依然成立。
\end{enumerate}
其中 (1) 是 $\Res\dashv\Ind$ 与 $\phi_{!}\dashv\phi^{*}$ 的形式推论：
\[
\Hom_L(\widetilde J\sigma_G,\sigma_L)\cong
\Hom_P\bigl(\Res^G_P(\sigma_G),\phi^{*}(\sigma_L)\bigr)\cong
\Hom_G(\sigma_G,\widetilde I\sigma_L).
\]
\end{remark}

\subsubsection{抛物诱导与 Jacquet 函子 II}

朴素的 Jacquet 函子与抛物诱导并不把酉表示送到酉表示（因为抛物子群一般不是幺模的）。
设 $\delta$ 是 $P$ 上的模特征标，即由
\[
(L_x)^{*}\mu^{r}_P=\delta(x)\mu^{r}_P
\]
给出，其中 $L_x$ 是 $x$ 的左平移，$\mu^{r}_P$ 是 $P$ 上（任一）右不变 Haar 测度。

\begin{example}
考虑 $B\subset\GL_2(K)$ 是上三角矩阵构成的标准 Borel 子群
$\bigl\{\begin{psmallmatrix}x&y\\&z\end{psmallmatrix}\bigr\}\subset\GL_2$，
则其左不变测度 $\mu^{l}$ 与右不变测度 $\mu^{r}$ 分别为
\[
d\mu^{l}=\frac{1}{|x|^{2}\cdot|z|}\,dx\,dy\,dz,
\qquad
d\mu^{r}=\frac{1}{|x|\cdot|z|^{2}}\,dx\,dy\,dz .
\]
按定义 $d\mu^{r}(\gamma g)=\delta(\gamma)d\mu^{r}(g)$，于是容易算出
\begin{multline*}
d\mu^{r}\Bigl(\begin{pmatrix}a&b\\&d\end{pmatrix}
\begin{pmatrix}x&y\\&z\end{pmatrix}\Bigr)
=\frac{d(ax)\wedge d(ay+bz)\wedge d(dz)}{|ax|\cdot|dz|^{2}}\\
=\frac{|a|}{|d|}\cdot d\mu^{r}\Bigl(\begin{pmatrix}x&y\\&z\end{pmatrix}\Bigr),
\end{multline*}
故 $\delta\bigl(\begin{psmallmatrix}a&b\\&d\end{psmallmatrix}\bigr)=\bigl|a\cdot d^{-1}\bigr|$。

更一般地，对 $L_{\underline n}\subset P_{\underline n}$（把 Levi 看作子群），
有 $P_{\underline n}$ 上的模特征标 $\delta_{\underline n}$，而 $\delta_{\underline n}|_{L_{\underline n}}$
可由
\[
\delta_{\underline n}(x)=\bigl|\det\bigl(\ad_{U_{\underline n}}(x)\bigr)\bigr|,
\qquad \forall x\in L_{\underline n}
\]
描述。这里 $\ad_{U_{\underline n}}$ 记 $L_{\underline n}$ 通过共轭作用在 $U_{\underline n}$ 上所得的表示。
\end{example}

这段讨论导向正规化 Jacquet 函子 $J$ 与抛物诱导 $I$ 的定义。

\begin{definition}
\[
J^G_P:=\delta^{-\frac12}\cdot\widetilde J^G_P,
\qquad
I^G_P:=\widetilde I^G_P\cdot\delta^{\frac12}.
\]
\end{definition}

\begin{remark}
按定义，对 $L_{\underline n}=\GL_{\underline n}$ 的任一光滑表示 $(\pi_L,W)$，有
\begin{multline*}
I^G_{P_{\underline n}}(W)=
\Bigl\{f\colon G\to W\ \Big|\ f\ \text{光滑},\ 
f(uhg)=\delta_{\underline n}^{1/2}(h)\,\pi_L(h)f(g)\\
\forall u\in U_{\underline n},\ h\in\GL_{\underline n},\ g\in G\Bigr\}.
\end{multline*}
\end{remark}

\subsubsection{超尖点表示}

\begin{definition}
设 $(\pi,V)\in\Repadm_G$ 是不可约容许表示。若对任一真抛物子群 $P<G$ 都有
$J^G_P(\pi)=0$，则称 $(\pi,V)$ 是\emph{超尖点的}。
\end{definition}

超尖点表示构成容许表示的基本构件。

\begin{lemma}
设 $\pi\in\Repinf_G$ 不可约（从而容许），则或者 $\pi$ 超尖点，或者 $\pi$ 是
$I^G_P(\sigma)$ 的子表示，其中 $P<G$ 是（真）抛物子群，$\sigma$ 是 $L$ 的不可约超尖点表示。
\end{lemma}

\begin{proof}
设 $\pi$ 不超尖点，则存在极小抛物子群 $P<G$ 使 $J^G_P(\pi)\neq 0$；$J^G_P\pi$ 有限长，
取其不可约商 $\sigma$。由伴随性
\[
0\neq\Hom_L(J^G_P\pi,\sigma)=\Hom_G(\pi,I^G_P\sigma),
\]
即 $\pi$ 同构于 $I^G_P\sigma$ 的某个子表示。按构造可以证明 $\sigma$ 超尖点。
\end{proof}

下述超尖点表示的描述（我认为属于 Harish-Chandra）很有用。令
$G^{0}=\{g\in\GL_n(K):\det g\in\mathcal{O}_K^{\times}\}$，$Z=Z(G)$ 为 $G$ 的中心，
则 $ZG^{0}$ 在 $G$ 中有限指数（这对所有约化 $p$-进群都成立，对 $\GL_n$ 指数为 $n$）。
其想法是：通常可以先研究 $G$ 的表示在 $G^{0}$ 上的限制，再单独处理中心。
本书后文会给出这一想法的应用。

\begin{lemma}
设 $(\pi,V)$ 是 $G$ 的不可约容许表示，则下列等价：\en{8}
\begin{enumerate}[label=(\arabic*),leftmargin=2.2em]
\item $\pi$ 超尖点；
\item 每个矩阵系数 $C_{v,\lambda}\colon g\mapsto\lambda(\pi(g)v)$（$G\to\mathbb{C}$）
      在 $G^{0}$ 上紧支集；
\item 每个矩阵系数模中心 $Z\subset G$ 紧支集。
\end{enumerate}
此外，若再设 $\pi^{\vee}$ 也不可约，则 (1)--(3) 还等价于
\begin{enumerate}[label=(\arabic*),leftmargin=2.2em]\setcounter{enumi}{3}
\item 存在一个矩阵系数模中心紧支集。
\end{enumerate}
\end{lemma}

\begin{proof}[略证]
底本此处只写「Add sketch in second draft」，未给证明。
\end{proof}

下面这个引理对我们有用：

\begin{lemma}
设 $K_0$ 是 $\GL_n(K)$ 的紧开子群。则 $G^{0}$ 的不可约超尖点表示中，
含有 $K_0$-不动向量的同构类只有有限多个。
\end{lemma}

\begin{proof}
这是 Bernstein 哈佛讲义第 3 节命题 21 推论的改写，见 [1]。
\end{proof}

\begin{remark}
下述从 $L_{\underline n}$ 到 $G$ 的抛物诱导的描述，将是 Bernstein--Zelevinsky 分类的特例。
\begin{itemize}[leftmargin=1.6em]
\item 设 $(\pi_L,W)$ 是 $L_{\underline n}=\GL_{n_1}\times\cdots\times\GL_{n_r}$ 的容许表示，
      由
\[
\pi_{\underline n}=\pi_1\otimes\cdots\otimes\pi_r,\qquad
W=V_1\otimes\cdots\otimes V_r
\]
给出，其中每个 $(\pi_i,V_i)$ 是 $\GL_{n_i}$ 的不可约超尖点表示。则表示
\[
\pi_1\times\cdots\times\pi_r:=I^G_{P_{\underline n}}(\pi_L)
\]
是 $G$ 的容许表示（如上所述）。它不可约当且仅当对任意 $i\neq j$ 的 $\pi_i,\pi_j$，
$\pi_i\neq|\det|\cdot\pi_j$。
\item 在 $\GL_2$ 情形，这说的正是：$\chi_1\times\chi_2$ 不可约当且仅当
      $\chi_1/\chi_2\neq|\cdot|^{\pm 1}$——这些就是 $\GL_2$ 的主序列表示。
\end{itemize}
\end{remark}
